% Zdroj: PDF priklady-leto-2023.pdf, fyzická str. 2-3. Značka: otázka studijního zaměření (neznačená starší sada). Zařazeno: S3.
\section{Shlukování}
\szzotazka{léto 2023}{8}{Shlukování}

\noindent\textbf{Zadání.}\par
Uvažujme data zadaná v tabulce níže a zobrazená v obrázku. Budeme na ně chtít aplikovat shlukovací algoritmy.

\begin{center}
\begin{minipage}{0.3\linewidth}
\centering
\begin{tabular}{cc}
\toprule
$x_1$ & $x_2$ \\
\midrule
1 & 1 \\
2 & 1 \\
3 & 2 \\
4 & 3 \\
4 & 2 \\
\bottomrule
\end{tabular}
\end{minipage}\hfill
\begin{minipage}{0.65\linewidth}
\centering
\includegraphics[width=\linewidth]{shlukovani-leto2023.png}
\end{minipage}
\end{center}

\begin{enumerate}
  \item Nejprve chceme aplikovat algoritmus $k$-means pro $k=2$. Na počátku jsme náhodně zvolili jako středy shluků body
    $(1,1)$ a $(3,2)$. Jaké budou nové středy shluků po jedné iteraci algoritmu $k$-means?
  \item Jak se změní polohy středů v dalších iteracích tohoto algoritmu?
  \item Popište některý z algoritmů pro hierarchické shlukování. Jaký by byl jeho výsledek na uvedených datech?
\end{enumerate}

\medskip
\noindent\textbf{Řešení.} \textit{(V~PDF bez náčrtu řešení; vlastní vzorové řešení, iterace ověřena numericky.)}\par
\begin{enumerate}[nosep]
  \item \emph{Přiřazení} (porovnáme vzdálenosti k~oběma středům): $(1,1)\to\mu_1$, $(2,1)\to\mu_1$ (vzdálenost $1<\sqrt2$), $(3,2)\to\mu_2$, $(4,3)\to\mu_2$, $(4,2)\to\mu_2$. \emph{Přepočet} (průměry shluků): $\mu_1=(1{,}5,\,1)$ a~$\mu_2=\big(\tfrac{3+4+4}{3},\,\tfrac{2+3+2}{3}\big)=\big(\tfrac{11}{3},\,\tfrac{7}{3}\big)$.
  \item Opakované přiřazení k~novým středům dá totéž rozdělení, takže se středy už nemění -- algoritmus zkonvergoval (po jediné iteraci).
  \item \emph{Aglomerativní} hierarchické shlukování: začneme s~pěti jednobodovými shluky a~opakovaně slučujeme dvojici nejbližších (single linkage = vzdálenost nejbližších bodů). Vzdálenost $1$ mají tři dvojice: $(1,1)$--$(2,1)$, $(3,2)$--$(4,2)$ a~$(4,2)$--$(4,3)$; ve výšce $1$ tak vzniknou skupiny $\{(1,1),(2,1)\}$ a~$\{(3,2),(4,2),(4,3)\}$, které se nakonec spojí ve výšce $\sqrt2$ (vzdálenost $(2,1)$--$(3,2)$). Řez dendrogramu na dva shluky dá totéž rozdělení jako $k$-means: $\{(1,1),(2,1)\}$ a~$\{(3,2),(4,2),(4,3)\}$.
\end{enumerate}
